% !TeX root = ../papers/nc-p3.tex
\section{反対三次曲線との接触による長い分岐列の排除}
\label{sec:branch-cubic-contact}
\begin{theorem}\label{thm:branch-cubic-contact}
定理\ref{thm:branch-centers-opposite-phases}の三配置のうち、残る可能性は $(3,3)$ だけである。
このとき両方の列の次数は $1,3,1$ で、偶位置における順・逆中心はともにそこでの中央三次曲線に属する。
\end{theorem}
\begin{proof}
順中心 $p$ の最初の直線を $L$、その次の実際の三次曲線を $A$ とする。
直線の位置には順・逆いずれの中心もないので、
命題\ref{prop:isolated-branch-blocks}とその逆により最初から二番目、最後から二番目の像は
完全三次系で埋め込まれる。$S$ を $L$ の消滅空間とする。
純関係自由な $L_3$ 鉛筆の三関係は $S\otimes V$ にあり、偶位置の適切な座標 $y$ で、
その評価は二行三列のHankel行列 $(y_{i+j})_{0\le i\le1,\,0\le j\le2}$ である。
その二次小行列式の零点スキームが反対三次曲線 $K$ である。
$L$ の外で後点があれば、それは $L$ を中心とする直線束への射影に三次埋込みを合成した値である。

同じ奇位置で別の列または同じ列の内部の曲線 $D$ を取り、両隣が埋め込まれた三次曲線とする。
その一般前点は動くので $D\ne L$。後点への写像が双有理なので、正規化からの
二つの射影座標の共通零因子は $\deg D-1$ である。$D$ 上で $L$ に入る点の前点は必ず $p$ で、
前の三次埋込みによりそのようなパラメータは高々一つである。さらにその重複度は高々1である。
実際、評価階数3の開集合では、三本の独立な列を選んで前点の座標比を周囲の開集合の正則関数で書ける。
この関数は $L$ 上で一定なので、$p$ で消える各座標の引戻しは $L$ のイデアルに入る。
前の写像が不分岐なため、正規化上でそのうち一つは位数1で消える。従って $\deg D\le2$。

長さ5なら第三曲線の両隣は上記の三次曲線である。長さ7なら第三曲線はその奇位置の初回直線でなく、
第四曲線はHankelの二者択一によりやはり三次曲線である。従って長い列では $d_3\le2$ で、
$d_1=1,d_2=d_4=3$、奇位置に順中心がないことから命題\ref{prop:branch-corridor-degree}は
$b_2=11-3d_3\ge5$ を与える。逆向きにも最後から二番目の三次曲線の順欠損は少なくとも5である。

$A\ne K$ なら $A$ 上の逆欠損因子はスキームとして $A\cap K$ に含まれる。
実際 $L\to A$ は同型で、各 $y\in A$ に対する前点ファイバーは $L$ と一点で交わる。
よって評価を $S$ で割った行列の階数は各点でちょうど1である。$S$ 外の単元を一つ選ぶと、
Hankelの二列とこの列から作る三次小行列式がその単元とHankel二次小行列式の積になる。
逆向きでも同じ包含が成立する。

$(3,7)$ なら長さ7の列の第二・第六曲線、$(5,5)$ なら一つの列の第二曲線と他方の第四曲線を取る。
いずれも同じ偶位置の相異なる三次曲線である。成分の一致は一般前点をたどって列の初めを復元すると矛盾する。
後者の三次曲線に来る奇曲線は前者の初回直線でないので、それは前者に対する $K$ である。
逆方向に読むと前者も後者に対する反対三次曲線である。従って両者の交叉は、
逆中心と順中心という相異なる二点で各々長さ5以上となる。
一方、相異なる有理正規三次曲線二本の交叉の長さは高々5である。
長さ6以上なら第一曲線の三本の独立な二次方程式は、第二曲線の上では
同じ長さ6の因子で消える六次式なので高々一次元を張る。よって独立な二式が両曲線を含む。
非平面的既約曲線を含む二次式は既約で、二式の完全交叉は次数4に過ぎない。これで長い列は消えた。
$(3,3)$ では次数公式と逆公式から第二曲線で $f_2=b_2=5$ を得る。
\end{proof}

\begin{proposition}\label{prop:branch-three-three-frame}
残った配置で偶位置の三次曲線を $C(a,b)=(a^3,a^2b,ab^2,b^3)$、
順・逆中心をそれぞれ $C(1,0),C(0,1)$ とする基底が取れる。
各奇位置では初回の直線と終端の直線が
$L(a,b)=(a,b,0,0)$ と $M(a,b)=(0,a,b,0)$ となり、動く二辺はパラメータを保って
$L\to C\to M$ をたどる。二偶位置の独立二項変数を $x,t$、間の奇位置を $v,v'$ とし、
\[
 h(x,t;v)=x_1t_1v_0-x_1t_0v_1+x_0t_0v_2
\]
と置く。窓の全係数を偶位置で三次曲線に制限して書いた多項式は
\[
 hh'B+v_3 h'P+v'_3 hQ+v_3v'_3D,\qquad
 h=h(x,t;v),\quad h'=h(t,x;v'),
\]
となる。$B,P,Q,D$ の二重次数は順に $(1,1),(2,2),(2,2),(3,3)$ であり、表示は一意である。
特にこの38次元線型空間内で $B$ は非退化でなければならない。
\end{proposition}
\begin{proof}
両中心は異なるので偶基底を選べる。定理\ref{thm:branch-rank-eleven-floor}の証明により
二つの奇直線は一点で交わる。この点の後点は逆中心、前点は順中心である。
従って各直線で隣の三次曲線から移したパラメータは表示の共通点で整合し、
交点の代表ベクトルを合わせて同時に四次元の基底へ延ばせる。
三次埋込みの座標は三次式の基底なので制限でテンソルの情報は失われない。
$v$ の最初の三係数は $M(x)$ と $L(t)$ をともに消す。従って分数体上でその外積の倍数であり、
三係数に共通の既約因子がないので多項式倍である。反対位置でも同様で、独立な文字 $v_3,v'_3$ で
分ければ表示と一意性を得る。逆に各項は四つの動く辺と両端の定点・直線の条件を満たす。

各六次元関係空間から隣接奇位置の第四文字を除く射影は単射である。核の元は全て純テンソルだからである。
従って両奇位置で第四文字を除いた $hh'B$ は、隣接二切断でいずれも階数6である。
各切断の行・列はそれぞれ固定された六空間に入る。実際、十二次元の座標に動く辺の五係数と
定点の二式を課すと一式だけ重なり六次元になる。$B=\sum B_{ij}x_0^{1-i}x_1^it_0^{1-j}t_1^j$
として両切断を展開すると、この六空間の座標小行列式は符号を除き $(\det B)^3$ である。
添付計算は両方の六座標を指定し多項式恒等式を検算する。
\end{proof}
